\clearpage
\VerdierPhysicalPageAnchor{0182}%
\begingroup\small\setlength{\parskip}{0pt}\setlength{\emergencystretch}{1.5em}
\setlength{\abovedisplayskip}{4pt plus 2pt minus 2pt}%
\setlength{\belowdisplayskip}{4pt plus 2pt minus 2pt}%

\VerdierSemanticAnchor{ch3d-1-3-1-naturality-cont}
le diagramme ci-après est commutatif :

\VerdierSemanticAnchor{ch3d-1-3-1-naturality-square}
\begin{center}
\begin{tikzcd}[column sep=4.5em,row sep=2.4em]
G(Z) \arrow[r,"{\delta(S)}"] \arrow[d,"{G(w)}"']
  & G(X)[1] \arrow[d,"{G(u)[1]}"] \\
G(Z') \arrow[r,"{\delta(S^{\prime})}"'] & G(X')[1]
\end{tikzcd}
\end{center}

\VerdierSemanticAnchor{ch3d-1-3-1-distinguished-triangle}
\noindent b) Pour toute suite exacte de $\mathcal A$,
$S=0\longrightarrow X\xrightarrow{u}Y\xrightarrow{v}Z\longrightarrow0$,
le triangle
$(G(X),\allowbreak G(Y),\allowbreak G(Z),\allowbreak G(u),\allowbreak
G(v),\allowbreak\delta(S))$ est distingué.

\VerdierSemanticAnchor{ch3d-1-3-2}
\medskip\noindent\textbf{1.3.2.}\quad
Soient $\mathcal A$ une catégorie abélienne et $u$ un morphisme de
$\operatorname{comp}(\mathcal A)$. On désigne par $[u]$ l'image de
$\dot u$ dans $D(\mathcal A)$. Soit
$S=0\longrightarrow X\xrightarrow{u}Y\xrightarrow{v}Z\longrightarrow0$
une suite exacte de $\operatorname{comp}(\mathcal A)$. Les propriétés
d'adjonction du foncteur cône (chap. I, 3.1.3) définissent un morphisme
$\mu(S)\colon c(u)\longrightarrow Z$. Le diagramme ci-après est commutatif
(chap. I, 3.2.2.4) :

\VerdierSemanticAnchor{ch3d-1-3-2-cone-square}
\begin{center}
\begin{tikzcd}[column sep=4.2em,row sep=2.8em]
Y \arrow[r,"p(u)"] \arrow[dr,"v"'] & c(u) \arrow[d,"\mu(S)"] \\
& Z
\end{tikzcd}
\end{center}

\VerdierSemanticAnchor{ch3d-1-3-2-components}
Un calcul simple montre que les composants du morphisme $\mu(S)$ sont :
\begin{equation}
\begin{aligned}
\mu(S)^i&\colon Y^i\oplus X^{i+1}\longrightarrow Z^i,\\
\mu(S)^i&=(v^i,0).
\end{aligned}
\tag{1.3.2.1}
\end{equation}
On en déduit que le morphisme $\mu(S)\colon c(u)\longrightarrow Z$ est un
épimorphisme et qu'on a une suite exacte :

\VerdierSemanticAnchor{ch3d-1-3-2-short-exact}
\[
0\longrightarrow c(\operatorname{id}_X)\longrightarrow c(u)
\longrightarrow Z\longrightarrow0.
\]
Introduisons alors le complexe $c(\mu(S))$ et soit
$m\colon c(\operatorname{id}_X)\longrightarrow c(\mu(S))[-1]$ le
morphisme de composants :

\VerdierSemanticAnchor{ch3d-1-3-2-morphism-m}
\begin{equation}
\begin{aligned}
m^i&\colon X^i\oplus X^{i+1}
  \longrightarrow Z^{i-1}\oplus Y^i\oplus X^{i+1},\\[2pt]
m^i&=
\begin{pmatrix}
0&0\\
u^i&0\\
0&\operatorname{id}
\end{pmatrix}.
\end{aligned}
\tag{1.3.2.2}
\end{equation}
\endgroup
